% !TEX root = ../Main.tex
\chapter{倒序相加法}

\section{\texorpdfstring{高斯怎么算 $1+2+\cdots+99$}{高斯怎么算 1+2+...+99}}

高斯小时候被老师罚算 $S=1+2+3+\cdots+99$。他的办法是\textbf{从两端往里配对}：
\[
S=1+2+3+\cdots+97+98+99,\qquad
\underbrace{1+99}_{100},\ \underbrace{2+98}_{100},\ \dots
\]
每一组的和都是 $100$。从 $1$ 到 $99$ 配对：
\[
\underbrace{1,2,\dots,49}_{\text{49 个}},\quad \boxed{50},\quad \underbrace{51,52,\dots,99}_{\text{49 个}},
\]
正好配成 $49$ 组、每组 $100$，正中间剩一个 $50$ 落单。所以
\[
S=49\times 100+50=4950.
\]

\section{奇偶讨论的麻烦}

再看 $S'=1+2+\cdots+99+100$。这回每组配成 $101$：
\[
\underbrace{1+100}_{101},\ \underbrace{2+99}_{101},\ \dots,\ \underbrace{50+51}_{101},
\]
\[
\underbrace{1,2,\dots,50}_{\text{50 个}}\,\Big|\,\underbrace{51,\dots,100}_{\text{50 个}},
\]
恰好分成 $50$ 组、没有落单，于是
\[
S'=50\times 101=5050.
\]

对比两次可以发现：\textbf{偶数个}这样的数相加，恰好两两均分；\textbf{奇数个}相加，中间那个数会落单，得单独处理。一个求和方法，竟要分奇偶两种情形——能不能避开？

\section{倒序相加：让每个数在"两式之间"找伴}

落单是因为我们在\textbf{同一个式子内部}给每个数找搭档。换个思路：写下两份同样的和，一份正着、一份倒着，让每个数到\textbf{另一份式子}里去找伴：
\[
\begin{array}{rcccccccl}
S=& 1 &+& 2 &+& \cdots &+& 99,\\
S=& 99 &+& 98 &+& \cdots &+& 1.
\end{array}
\]
上下对齐相加，\textbf{每一列都是 $1+99=100$}，一共 $99$ 列，所以
\[
2S=100\times 99\quad\Longrightarrow\quad S=\frac{100\times 99}{2}=50\times 99=4950.
\]
这下不论项数是奇是偶，每一列都成对，干干净净，再无奇偶之分。

\state{倒序相加（高斯求和）}{
    把 $S=1+2+\cdots+n$ 倒着再写一份对齐相加，每列都是 $n+1$，共 $n$ 列：
    \[
    2S=n(n+1)\quad\Longrightarrow\quad S=1+2+\cdots+n=\frac{n(n+1)}{2}.
    \]
}
